\chapter{Introdução}
\label{cap:introducao}

As regras que um aluno de mestrado ou doutorado precisa cumprir estão espalhadas por
documentos de naturezas diferentes. No MDCC/UFC, o prazo de vínculo vem da resolução
que fixa as normas gerais da pós-graduação da universidade \cite{UFC2015}; os créditos,
a composição das bancas e a antecedência da divulgação de uma defesa vêm do regimento
do programa \cite{MDCC2016}; os exames de proficiência aceitos e o prazo para pedir
aproveitamento de disciplinas aparecem apenas em páginas de procedimentos do site. Uma
pergunta simples, como quantos créditos faltam para concluir o curso, pode exigir a
leitura de mais de um documento, e as respostas nem sempre coincidem: o regimento fala
em 24 créditos em disciplinas, enquanto a resolução fala em um mínimo de 30 créditos,
dos quais seis correspondem à dissertação.

Modelos de linguagem de grande escala (LLMs) respondem bem a perguntas em linguagem
natural, mas, sem acesso aos documentos, tendem a produzir regras plausíveis e
inexistentes \cite{Ji2023}. Em um domínio normativo esse erro tem custo direto para o
aluno. A geração aumentada por recuperação (RAG) reduz o problema ao entregar ao modelo
os trechos relevantes no momento da pergunta \cite{Lewis2020,Gao2023}, e permite que
cada afirmação seja rastreada até a sua fonte.

Este relatório descreve a construção e a avaliação de um assistente RAG para as normas do MDCC. As decisões
de projeto seguem cinco perguntas: qual modelo usar e por quê; qual estratégia de
adaptação (\textit{prompting}, RAG ou ajuste fino); como avaliar de forma automática e
humana; quanto custa cada resposta em tokens, latência e memória; e que riscos de
viés, privacidade e segurança o sistema traz. As contribuições são:
(i) um corpus versionado das normas do programa, com uma divisão em trechos que
preserva artigos e itens, de modo que a citação vem dos metadados e não da memória do
modelo; (ii) um conjunto de perguntas de avaliação em que cada resposta de referência
tem um trecho literal do documento como evidência; e (iii) uma comparação de três
geradores com métricas de recuperação, de qualidade da resposta, de custo e de
segurança, todas reproduzíveis a partir de arquivos.

O Capítulo~\ref{cap:trabalhos-relacionados} apresenta os trabalhos relacionados. O
Capítulo~\ref{cap:metodologia} descreve o corpus, o sistema e o protocolo de avaliação. O
Capítulo~\ref{cap:resultados} traz os resultados, discutidos no
Capítulo~\ref{cap:discussao}. O Capítulo~\ref{cap:conclusao} conclui.
